% ============================================================================
%  CONCLUSION GÉNÉRALE ET PERSPECTIVES
% ============================================================================
\chapter*{Conclusion générale et perspectives}
\addcontentsline{toc}{chapter}{Conclusion générale et perspectives}
\markboth{Conclusion générale et perspectives}{Conclusion générale et perspectives}

Trois manières de franchir un canal bruité simulé ont été comparées sous un budget radio identique : la
chaîne séparée simulée bit à bit, un codage conjoint appris sur une perte pixel, et le même réseau entraîné
avec une perte sémantique. Le dispositif a été construit pour que la comparaison soit recevable —
architecture identique entre les deux systèmes appris, puissance normalisée, canal et tirages communs, chaîne
numérique non approchée, juges absents de toute fonction de perte.

\paragraph{Résultats.} Ajouter à la perte un terme de caractéristiques et un terme de tâche, et rien d'autre,
échange environ \SI{0.5}{\decibel} de PSNR contre 10 à 14~points de précision chez un juge indépendant, en
AWGN comme en Rayleigh. En supprimant les bits, le codage conjoint supprime l'effet falaise : à
\SI{0}{\decibel}, là où la chaîne séparée réelle tombe au niveau du hasard, les systèmes appris restent
exploitables. Enfin, un modèle entraîné à SNR fixe est inutilisable ailleurs : l'entraînement sur une plage
est la condition d'un système déployable.

\paragraph{Limites.} Le canal est simulé et la CSI supposée parfaite sous Rayleigh. Les modèles n'ont vu que
des images $32\times32$ : sur Kodak et en vidéo, la borne idéale de la chaîne séparée reste devant en PSNR,
ce qui est un résultat de transfert et non une défaite de principe. La borne de capacité, enfin, n'est pas un
système constructible.

\paragraph{Extensions.} Quand la tâche du récepteur est connue, la transmission orientée tâche atteint une
précision comparable avec 32 fois moins de bande. Quand le contenu est temporellement redondant, le
renouvellement conditionnel divise la bande par trois sur une scène statique — et doit être désactivé
ailleurs.

\paragraph{Perspectives.} À court terme, l'\textbf{entraînement sur images naturelles}, seul moyen de
trancher la comparaison avec JPEG sur Kodak, et le \textbf{pilotage automatique du renouvellement
conditionnel} par une mesure de mouvement. À moyen terme : une \textbf{estimation imparfaite du canal},
l'étude du gain sémantique pour \textbf{d'autres tâches} que la classification, et une \textbf{quantification
des symboles émis}, nécessaire pour rencontrer un émetteur radio réel. À plus long terme, la validation sur
une radio logicielle reste l'étape qui manque pour passer d'une démonstration en simulation à un résultat de
transmission.
